\documentclass[12pt]{article}
\usepackage{cmap}
\usepackage[T2A]{fontenc}
\usepackage[utf8]{inputenc}
\usepackage[english,russian]{babel}
\usepackage{geometry}
\usepackage{titlesec}
\usepackage{tocloft}
\usepackage{longtable}
\usepackage{array}
\usepackage{amsmath}
\setlength{\emergencystretch}{1em}
\geometry{a4paper, left=2.5cm, right=2.5cm, top=2.5cm, bottom=2.5cm}
\pagenumbering{gobble}

\begin{document}

Мне нужно было удалить из базы лишние поля и добавить в таблицу фактов агрегатные строки.

\section*{1. Какие поля убрал}

\noindent\textbf{Из \texttt{dim\_date}:}
\begin{itemize}
    \item \texttt{day\_no} - Конкретные дни излишни.
    \item \texttt{quarter\_no} - Районы лишние для аналитики задолженностей и оплат по городам.
    \item \texttt{month\_start} - Начало месяца есть в \texttt{billing\_month}.
\end{itemize}

\noindent\textbf{Из \texttt{dim\_apartment}:}
\begin{itemize}
    \item \texttt{source\_apartment\_id} - Дублирует \texttt{apartment\_key}.
    \item \texttt{passport\_date} - Не используется в аналитике. Не анализируем даты регистрации квартир
    \item \texttt{floor\_no} - Не изучаем данные по этажам.
    \item \texttt{street}, \texttt{house\_number}, \texttt{apartment\_number} - Не анализируем конкретные квартиры и адреса.
\end{itemize}

\noindent\textbf{Из \texttt{fact\_billing}:}
\begin{itemize}
    \item \texttt{due\_date} - Срок оплаты можно вычислить из расчётного месяца: 10-е число следующего месяца.
    \item \texttt{snapshot\_date} - Дата снимка одна для всей базы, 31 июля 2025 года.
\end{itemize}

\section*{2. Какие агрегированные строки добавляем}

Добавил такие агрегаты:

\begin{center}
\textbf{Год $\times$ город $\times$ услуга $\times$ интервал просрочки.}
\end{center}

Интервалов три: до 30 дней, 31--90 дней и более 90 дней.

Для каждой группы заранее вычисляем два показателя: общую сумму задолженности и количество записей с долгом.

Эти значения рассчитываем по исходным строкам один раз и сохраняем дополнительными строками в \texttt{fact\_billing}. Чтобы отличать их от обычных начислений, используем отрицательные \texttt{apartment\_key}.


\end{document}
